\begin{figure}[h]
  \centering
  \resizebox{\linewidth}{!}{%
  \begin{tikzpicture}[
      node distance=7mm and 9mm,
      font=\scriptsize,
      every node/.style={align=center},
      source/.style={draw, rounded corners=1pt, fill=blue!8,
                     minimum width=18mm, minimum height=5mm, inner sep=1pt},
      stage/.style={draw, thick, rounded corners=2pt, fill=gray!10,
                    minimum width=26mm, minimum height=10mm, inner sep=3pt},
      metric/.style={draw, rounded corners=1pt, fill=orange!12,
                     minimum width=26mm, minimum height=5mm, inner sep=1pt},
      output/.style={draw, thick, rounded corners=2pt, fill=green!10,
                     minimum width=30mm, minimum height=10mm, inner sep=3pt},
      flow/.style={-{Latex[length=2mm]}, thick},
    ]
    % Column 1: sources.
    \node[source] (s1) {NeurIPS 1987--2025 \\ {\footnotesize ICLR 2013+, ICML 2008+}};
    \node[source, below=of s1] (s2) {arXiv 1991--2025 \\ {\footnotesize (ML and non-ML)}};
    \node[source, below=of s2] (s3) {PubMed 1990--2025};
    \node[above=1mm of s1, font=\scriptsize\bfseries] {Data sources};

    % Column 2: canonical schema.
    \node[stage, right=12mm of s2.east] (schema)
      {Canonical schema \\ (Pydantic validated) \\
       \texttt{paper\_id, venue, year,} \\
       \texttt{title, abstract, authors,} \\
       \texttt{arxiv\_primary\_category}};

    % Column 3: per paper metrics.
    \node[metric, right=14mm of schema.north east, yshift=-2mm] (m1)
      {Classical readability \\ (15 scores, \texttt{textstat})};
    \node[metric, below=2mm of m1] (m2)
      {Hohmann writing style \\ (9 metrics, spaCy)};
    \node[metric, below=2mm of m2] (m3)
      {Acronym density + reuse \\ (Barnett 2020 algorithm)};
    \node[metric, below=2mm of m3] (m4)
      {LLM as judge readability \\ (Cachola 2025 prompts)};
    \node[metric, below=2mm of m4] (m5)
      {Citations (OpenAlex, S2)};
    \node[above=1mm of m1, font=\scriptsize\bfseries] {Per paper metrics};

    % Column 4: aggregation.
    \node[stage, right=12mm of m3.east] (agg)
      {Aggregation \\ per venue-year \\ per arXiv category};

    % Column 5: outputs.
    \node[output, right=11mm of agg.north east, yshift=-5mm] (ev)
      {\textbf{Evidence} \\ volume, readability, \\ acronyms, fragmentation};
    \node[output, below=3mm of ev] (pol)
      {\textbf{Recommendations} \\ seven measurable standards};

    % Flow arrows: sources to schema.
    \foreach \src in {s1, s2, s3}
      \draw[flow] (\src.east) -- (schema.west);

    % Schema to metrics.
    \foreach \met in {m1, m2, m3, m4, m5}
      \draw[flow] (schema.east) -- (\met.west);

    % Metrics to aggregation.
    \foreach \met in {m1, m2, m3, m4, m5}
      \draw[flow] (\met.east) -- (agg.west);

    % Aggregation to outputs.
    \draw[flow] (agg.east) -- (ev.west);
    \draw[flow] (agg.east) -- (pol.west);
  \end{tikzpicture}%
  }
  \caption{\textbf{End to end workflow.} Three data sources feed a
  canonical Pydantic-validated schema. Five per-paper metric
  families are computed \revfour{on abstracts, and acronym density
  also on titles,} then
  aggregated per venue-year and per arXiv primary category.
  The outputs feed the empirical evidence
  (Sections~\ref{sec:readability} through~\ref{sec:volume}) and
  the recommendations
  (Section~\ref{sec:policy}). Pipeline details appear in
  Appendix~\ref{app:workflow}.}
  \label{fig:workflow}
\end{figure}
